%% LyX 2.4.4 created this file.  For more info, see https://www.lyx.org/.
%% Do not edit unless you really know what you are doing.
\documentclass[english]{article}
\usepackage[T1]{fontenc}
\usepackage[latin9]{inputenc}
\usepackage{enumitem}
\usepackage{amsmath}
\usepackage{amsthm}
\usepackage{geometry}
\geometry{verbose,tmargin=1in,bmargin=1in,lmargin=1in,rmargin=1in}

\makeatletter
%%%%%%%%%%%%%%%%%%%%%%%%%%%%%% Textclass specific LaTeX commands.
\numberwithin{equation}{section}
\numberwithin{figure}{section}
\newlength{\lyxlabelwidth}      % auxiliary length 

%%%%%%%%%%%%%%%%%%%%%%%%%%%%%% User specified LaTeX commands.
\usepackage{fancyhdr}
\usepackage{lastpage}
\pagestyle{fancy}
\usepackage{enumitem}
\usepackage{ifthen}
\everymath=\expandafter{\the\everymath\displaystyle}
%\DeclareMathSizes{10}{10}{10}{10}
% Added by lyx2lyx
\setlength{\parskip}{\smallskipamount}
\setlength{\parindent}{0pt}

\makeatother

\usepackage{babel}
\begin{document}
\renewcommand{\author}{Dr.\,James Ashe}

\newcommand{\classNum}{331}

\newcommand{\classSec}{A}

\newcommand{\examType}{EXAM}

\renewcommand{\date}{October 12, 2012}

%%First page's header%%%%%%%%%%%%%%%%%%%%%%
\fancypagestyle{firststyle} %%header settings for first pagr
{
	\fancyhead[L]{MTH \classNum{}(\classSec{}) \author{}\\
	\textbf{\examType{}} \date{}}
	\fancyhead[C]{\textbf{Name:}}
	\setlength{\headsep}{14pt}
}
%%%%%%%%%%%%%%%%%%%%%%%%%%%%%%%%%%%%%%%%%%

%%Next page's header%%%%%%%%%%%%%%%%%%%%%%%
\setlist{nolistsep}
\lhead{\classNum{}(\classSec{})} 
\chead{\textbf{Name:}} 
\cfoot{\thepage{}~of~\pageref{LastPage}} 
\setlength{\headsep}{5pt} 
%%%%%%%%%%%%%%%%%%%%%%%%%%%%%%%%%%%%%%%%%%%

%%Various settings; see the preamble for other settings%%%%%
%%\setenumerate[0]{leftmargin=12pt,itemindent=0pt,labelwidth=0pt}
\renewcommand{\labelenumi}{\textbf{\arabic{enumi}.}}
\renewcommand{\labelenumii}{\textbf{\alph{enumii}.}}
\thispagestyle{firststyle}
%%%%%%%%%%%%%%%%%%%%%%%%%%%%%%%%%%%%%%%%%%

\textbf{Justify your answers and show your work. All questions are
of equal value.}
\begin{enumerate}[resume]
\item Find an explicit formula for the general $n^{\mbox{th}}$ term of
the sequence, $\left\{ a_{n}\right\} _{n=0}^{\infty}=\{1,2,4,8,16,\dots\}$.
\vspace{1in}
\item Write the infinite series, $0.3+0.03+0.003+\cdots$, using sigma notation
(i.e. using $\Sigma$) . \vspace{1in}
\item Write the first 3 terms of the sequence of partial sums of the series
${\displaystyle \sum_{k=1}^{\infty}\frac{3^{k}}{2}}$.\vfill
\end{enumerate}
\emph{Evaluate the geometric series or state that it diverges. }
\begin{enumerate}[resume]
\item ${\displaystyle \sum_{k=0}^{\infty}\frac{2^{k}}{7^{k}}}$.\vfill
\item ${\displaystyle \sum_{m=2}^{\infty}\left(\frac{1}{3}\right)^{m}}$.\vfill\newpage{}
\end{enumerate}
\emph{Find the limit of the following sequences or determine that
the limit does not exist. }
\begin{enumerate}[resume]
\item ${\displaystyle \lim_{n\rightarrow\infty}\left\{ \frac{n^{4}-1}{n^{4}+10}\right\} }$.\vfill
\item ${\displaystyle \lim_{n\rightarrow\infty}\left\{ \frac{n^{0.01}}{\ln n}\right\} }$.\vfill
\end{enumerate}
\emph{Use the Divergence Test to determine whether the following series
diverge, or state that the test is inconclusive (use your results
from above).}
\begin{enumerate}[resume]
\item ${\displaystyle \sum_{k=1}^{\infty}\frac{k^{4}-1}{k^{4}+10}}$.\vfill
\item ${\displaystyle \sum_{k=1}^{\infty}\frac{k^{0.01}}{\ln k}}$.\vfill\newpage{}
\end{enumerate}
\emph{Using the Squeeze Theorem, find the limit of the following sequence
or determine that the limit does not exist.}
\begin{enumerate}[resume]
\item ${\displaystyle \lim_{n\rightarrow\infty}\left\{ \frac{\sin^{2}n}{2^{n}}\right\} }$.\vfill
\end{enumerate}
\emph{Use the Integral Test to determine whether the series converge,
diverge, or state why the test cannot be used. Check that the conditions
of the test are satisfied.}
\begin{enumerate}[resume]
\item ${\displaystyle \sum_{k=1}^{\infty}\frac{k}{e^{k}}}$. For $f(x)=\frac{x}{e^{x}}$,
$f'(x)=\left(1-x\right)e^{-1}$ and $\int f(x)\,dx=2xe^{-1}$.\vfill
\item ${\displaystyle \sum_{k=2}^{\infty}\frac{1}{(2x-3)^{2}}}$. For $f(x)=\frac{1}{(2x-3)^{2}}$,
$f'(x)=\frac{-4}{(2x-3)^{2}}$ and $\int f(x)\,dx=\frac{-1}{2(2x-3)}$.\vfill
\end{enumerate}
\emph{Use the Ratio Test to determine whether the following series
converges.}
\begin{enumerate}[resume]
\item ${\displaystyle \sum_{k=1}^{\infty}\frac{3}{k!}}$.\vfill\newpage{}
\end{enumerate}
\emph{Use the Root Test to determine whether the following series
converges. }
\begin{enumerate}[resume]
\item ${\displaystyle \sum_{k=1}^{\infty}\left(\frac{k-1}{3k}\right)^{k}}$.\vfill
\end{enumerate}
\emph{Use the Comparison Test or Limit Comparison Test to determine
whether the following series converge. HINT: ${\displaystyle \sum\frac{1}{k^{p}}}$
converges for $p>1$ and diverges for $p\leq1$.}
\begin{enumerate}[resume]
\item ${\displaystyle \sum_{k=1}^{\infty}\frac{1}{k^{2}+4}}$.\vfill
\end{enumerate}
\emph{BONUS: Use the formal definition of the limit of a sequence
to prove that ${\displaystyle \lim_{n\rightarrow\infty}\left(3+\frac{1}{n}\right)}=3$.\vfill}
\end{document}
